\documentclass[11pt,a4paper]{article}
\usepackage[T1]{fontenc}
\usepackage[utf8]{inputenc}
\usepackage[czech]{babel}
\usepackage{lmodern,microtype}
\usepackage{amsmath,amssymb,mathtools}
\usepackage{geometry}
\geometry{left=19mm,right=19mm,top=21mm,bottom=21mm}
\usepackage[dvipsnames]{xcolor}
\usepackage{enumitem,tabularx,booktabs}
\usepackage[most]{tcolorbox}
\usepackage{fancyhdr,lastpage}
\usepackage[hidelinks]{hyperref}
\definecolor{deep}{HTML}{203955}
\definecolor{teal}{HTML}{176E76}
\definecolor{slate}{HTML}{506173}
\definecolor{pale}{HTML}{F3F7F9}
\setlength{\parindent}{0pt}
\setlength{\parskip}{5pt}
\setlist[enumerate]{leftmargin=6mm,itemsep=5pt,topsep=3pt}
\pagestyle{fancy}\fancyhf{}
\fancyhead[L]{\sffamily\footnotesize\color{slate} MAT · 4. ročník · kombinatorika}
\fancyhead[R]{\sffamily\footnotesize\color{slate} Hodina \LESSONNUM\, / 112}
\fancyfoot[L]{\sffamily\footnotesize\color{slate} Výukový podklad · samostatný dokument}
\fancyfoot[R]{\sffamily\footnotesize\color{slate} \thepage/\pageref{LastPage}}
\renewcommand{\headrulewidth}{0.35pt}
\renewcommand{\footrulewidth}{0pt}
\newcommand{\sect}[1]{\par\vspace{2mm}{\Large\sffamily\bfseries\color{deep} #1}\par\vspace{-0.5mm}\noindent\color{teal}\rule{\linewidth}{0.7pt}\color{black}\vspace{1mm}}
\tcbset{theobox/.style={enhanced,breakable,colback=pale,colframe=teal!45!white,boxrule=.6pt,arc=1mm,left=3mm,right=3mm,top=2mm,bottom=2mm}}
\newcommand{\LESSONNUM}{38}
\begin{document}

{\Huge\sffamily\bfseries\color{deep} Pascalův trojúhelník a počítání cest}\par
{\sffamily\large\color{teal} Výukový podklad · 38. vyučovací hodina}\par
\textbf{Hodinová dotace:} 45 minut\quad \textbf{Ročník:} 4. SPŠ\quad \textbf{Rozpis:} TP 2.3\par
\begin{tcolorbox}[theobox]\textbf{Výukový cíl.} Sestavit Pascalův trojúhelník, určit jeho řádky pomocí kombinačních čísel a využít ho při počítání mřížkových cest.\end{tcolorbox}
\sect{Teorie a vysvětlení}
Pascalův trojúhelník vzniká tak, že na krajích píšeme \(1\) a každý vnitřní člen je součtem dvou členů bezprostředně nad ním. \textbf{Řádek indexovaný \(n\)} obsahuje \(\binom n0,\binom n1,\ldots,\binom nn\). Začínáme \(n=0\), první řádek tedy obsahuje jedinou jedničku.\par
\[\begin{array}{c}1\\1\quad1\\1\quad2\quad1\\1\quad3\quad3\quad1\\1\quad4\quad6\quad4\quad1\\1\quad5\quad10\quad10\quad5\quad1\end{array}\]
\textbf{Cesty v mřížce:} Na nejkratší cestě složené pouze z kroků doprava \(D\) a nahoru \(N\) určuje každý výsledek pořadí kroků. Pro \(a\) kroků doprava a \(b\) nahoru je cest \(\binom{a+b}{a}=\binom{a+b}{b}\).
\sect{Řešené příklady}
\textbf{Příklad 1.} Doplň řádek \(n=6\) Pascalova trojúhelníku.\par
\textit{Řešení.} Každé nové číslo je součtem dvou předcházejících nad ním: \(1,6,15,20,15,6,1\).\par\smallskip
\textbf{Příklad 2.} Urči součet čísel v řádku \(n=5\).\par
\textit{Řešení.} Součet \(1+5+10+10+5+1=32\). Také přímo \(2^5=32\), protože řádek tvoří \(\binom5k\).\par\smallskip
\textbf{Příklad 3.} Kolik je nejkratších cest z \((0,0)\) do \((3,4)\), pouze doprava a nahoru?\par
\textit{Řešení.} Každá cesta je posloupnost sedmi kroků, v níž jsou právě tři \(D\) a čtyři \(N\). Výběr tří pozic pro \(D\): \(\binom73=35\).\par\smallskip
\textbf{Příklad 4.} Kolik takových cest z \((0,0)\) do \((4,3)\) prochází bodem \((2,1)\)?\par
\textit{Řešení.} K bodu \((2,1)\): dvě \(D\), jedno \(N\), tedy \(\binom32=3\) cesty. Odtud do \((4,3)\): dvě \(D\) a dvě \(N\), tedy \(\binom42=6\) cest. Pravidlem součinu \(3\cdot6=18\).\par\smallskip
\textbf{Příklad 5.} Najdi v řádku \(n=8\) číslo se spodním indexem \(k=2\).\par
\textit{Řešení.} Je to \(\binom82=28\). Můžeme ho získat Pascalovým součtem \(\binom71+\binom72=7+21=28\).\par\smallskip
\clearpage\sect{Procvičování pro studenty}
\begin{enumerate}
\item Vypiš řádek \(n=4\).
\item Urči střední číslo řádku \(n=6\).
\item Urči součet řádku \(n=7\).
\item Kolik nejkratších cest vede z \((0,0)\) do \((2,3)\)?
\item Kolik nejkratších cest vede z \((0,0)\) do \((4,2)\) přes \((2,1)\)?
\end{enumerate}
\sect{Řešení a zdůvodnění k procvičování}
\begin{enumerate}
\item \(1,4,6,4,1\).
\item \(\binom63=20\).
\item \(2^7=128\).
\item Pět kroků, z toho dva doprava: \(\binom52=10\).
\item První část má \(\binom32=3\) cesty, druhá také \(\binom32=3\), celkem \(9\).
\end{enumerate}
\sect{Metodické poznámky}
\textbf{Typické chyby:}
\begin{itemize}[leftmargin=5mm,itemsep=1pt]
\item Počáteční řádek indexujeme \(n=0\); „sedmý řádek“ bez upřesnění indexace může být nejednoznačný.
\item V mřížce počítáme pouze nejkratší cesty, tedy bez kroků zpět.
\item Při cestě přes daný bod se počty obou úseků násobí, ne sčítají.
\end{itemize}
\textbf{Doporučený průběh 45 minut:} 0–10 sestavení trojúhelníku; 10–20 indexy a součty; 20–32 mřížkové cesty; 32–42 procvičení; 42–45 rekapitulace.\par
\textbf{Návaznost:} Koeficienty Pascalova trojúhelníku se v další hodině stanou koeficienty při rozvoji mocnin dvojčlenu.\par
\vspace{2mm}\noindent{\footnotesize\color{slate} \textbf{Východiska:} rozpis výuky MAT, 4. ročník, hodiny 29--41; tematický plán SPŠ; srovnání tematického rozsahu a náročnosti s realisticky.cz, oddíl Kombinatorika (lekce 090101--090119). Výklad, příklady a řešení jsou nově autorsky sestavené, nikoli převzatá sbírka.}\par
\end{document}